\section{Thiết kế Dialogue Management}

Dialogue Management là lớp quyết định chatbot sẽ làm gì sau khi NLU đã phân tích câu hỏi. Đây là phần biến intent thành action.

Công thức khái quát có thể mô tả như sau:

\[
action_t = Policy(intent_t, confidence_t, entity_t, slots_t, history_t)
\]

Trong đó, \texttt{intent\_t} là ý định hiện tại của người dùng, \texttt{confidence\_t} là độ tin cậy của intent, \texttt{entity\_t} là thực thể được trích xuất, \texttt{slots\_t} là trạng thái hội thoại đang lưu và \texttt{history\_t} là lịch sử các lượt hội thoại trước.

Rasa sử dụng các thành phần như RulePolicy, MemoizationPolicy hoặc TEDPolicy để quyết định action. Với các luồng đơn giản, RulePolicy được dùng để ánh xạ trực tiếp intent sang action, ví dụ \texttt{chao\_hoi} sang \texttt{utter\_chao\_hoi}, \texttt{cam\_on} sang \texttt{utter\_cam\_on}, \texttt{tam\_biet} sang \texttt{utter\_tam\_biet}.

\begin{figure}[H]
    \centering
    \includegraphics[width=0.95\textwidth]{\detokenize{content/source/rasa/Quan hệ giữa DM, domain, NLG template và action server.png}}
    \caption{Quan hệ giữa Dialogue Manager, domain, NLG template và action server}
    \label{fig:rasa-dm-domain-nlg-action}
\end{figure}

Hình~\ref{fig:rasa-dm-domain-nlg-action} thể hiện mối quan hệ giữa intent ranking, tracker history, slot state, Dialogue Manager, template response và custom action. Nếu action được chọn là template, hệ thống lấy phản hồi từ domain. Nếu action cần xử lý nghiệp vụ, Rasa gọi Action Server.

Với các luồng cần ngữ cảnh, tracker và slot được sử dụng để ghi nhớ thông tin. Ví dụ, người dùng hỏi ``Ngành An toàn thông tin xét tổ hợp nào?'' rồi hỏi tiếp ``Thế điểm chuẩn năm ngoái là bao nhiêu?''. Ở lượt thứ hai, câu hỏi không nhắc lại tên ngành, nhưng hệ thống có thể dùng slot \texttt{nganh = An toàn thông tin} từ lượt trước để hiểu ngữ cảnh.

Thiết kế Dialogue Management cần bảo đảm chatbot xử lý được ba nhóm luồng: luồng một lượt, luồng nhiều lượt và luồng fallback. Luồng một lượt xảy ra khi người dùng hỏi rõ và hệ thống trả lời ngay. Luồng nhiều lượt xảy ra khi người dùng hỏi thiếu thông tin, bot hỏi lại và người dùng bổ sung. Luồng fallback xảy ra khi câu hỏi mơ hồ hoặc ngoài phạm vi, bot xin làm rõ hoặc ghi nhận câu hỏi.

